\chapter{Excavation and Groundworks Risk Assessment}
\label{chap:Excavation and Groundworks Risk Assessment}

%---------------------------- SECTION ----------------------------
\section{Why this assessment matters in construction}

Excavation work accounts for a disproportionate share of construction fatalities in the United Kingdom. Collapses of unsupported trench sides are the single most common cause of excavation-related death, and the physics are unforgiving: one cubic metre of soil weighs approximately 1.5 tonnes, and a person buried to the chest by a trench collapse will typically die from compressive asphyxiation within minutes, long before rescue equipment can be mobilised. The HSE's analysis of fatal accidents in excavations consistently identifies a failure to support trench sides, a failure to batter back excavation faces, and a failure to provide safe means of access and egress as the primary causal factors. None of these failures are the result of unknown hazards or unavailable controls; they are the result of inadequate risk assessment, inadequate supervision, and cost or programme pressure overriding safe working practice.

The hazard environment in excavation and groundworks extends well beyond the collapse of unsupported sides. Buried services --- gas mains, electricity cables, water mains, telecommunications ducts, and drainage infrastructure --- represent a persistent and severe risk on any excavation in an existing urban or suburban environment. Between 2016 and 2022, the HSE recorded an average of over fifty serious injuries per year attributable to contact with buried services during excavation work. Striking a live gas main with a mechanical excavator can cause an immediate explosion; striking a high-voltage electricity cable can cause electrocution of the plant operator and operatives in the vicinity. The pre-excavation survey, the CAT and Genny scan, and the verification of service drawings before excavation commences are not optional refinements to good practice --- they are fundamental risk controls that must appear in every groundworks risk assessment.

Groundwater and flooding introduce a further dimension of risk that is frequently underestimated at the assessment stage. An excavation that is dry at the time of the initial risk assessment may flood within hours of heavy rainfall, destabilising previously adequate shoring, undermining adjacent foundations, and creating a confined-space atmosphere risk from oxygen depletion and carbon dioxide accumulation in flooded voids. The assessment must specifically address the groundwater regime at the site, the drainage management plan during excavation, and the emergency procedures in the event of unplanned flooding.

The interaction between deep excavations and adjacent structures represents one of the most technically demanding aspects of the groundworks risk assessment. Underpinning of adjacent buildings, monitoring of adjacent structures during excavation, the assessment of surcharge loads from plant and spoil heaps adjacent to excavation edges, and the design of temporary retaining structures all require input from competent structural engineers. The groundworks risk assessment must define at what point engineering judgement is required and ensure that this judgement is obtained before excavation commences, not after a problem has manifested.

\barquote{``An unshored trench is not a calculated risk --- it is an act of negligence disguised as a time-saving measure.''}

\example{Site Reality}{A groundworks contractor is excavating a 2.8-metre deep trench for a new foul sewer connection in a residential street. The soil is a mixture of made ground and London Clay. The contractor has not obtained updated service drawings from the relevant undertakers, relying instead on a set of drawings dated 2009. A mechanical excavator strikes a live 11\,kV electricity cable at 1.4 metres depth that does not appear on the 2009 drawings; the cable was diverted in 2015 as part of a road widening scheme. The excavator operator receives a flashover injury. The trench sides are not shored or battered back, and the investigation reveals that the spoil heap has been deposited within 0.8 metres of the trench edge, adding surcharge loading to an already unstable trench profile. The risk assessment made no reference to service verification, no reference to shoring requirements, and no reference to the proximity restrictions for spoil heaps.}

%---------------------------- SECTION ----------------------------
\section{Legal and professional framework}

The primary duty framework for excavation work in construction is established by the Construction (Design and Management) Regulations 2015, which require that excavations are identified in the pre-construction information where known hazards exist and that the Construction Phase Plan addresses the specific management arrangements for excavation. The Health and Safety at Work etc.\ Act 1974 imposes the general duty to ensure, so far as is reasonably practicable, the health and safety of employees and others who may be affected by excavation activities. The Management of Health and Safety at Work Regulations 1999 require a suitable and sufficient risk assessment of all excavation activities before work commences.

HSE guidance HSG47 (Avoiding Danger from Underground Services) provides the industry-standard framework for the management of buried service risk. The Work in Compressed Air Regulations 1996 apply where compressed-air working is used in tunnelling or caisson operations. BS 6031:2009 Code of Practice for Earthworks provides technical guidance on the design and execution of earthworks and excavations that informs the competency standard for groundworks contractors. PUWER 1998 applies to all excavation plant and shoring equipment, requiring that it is suitable, maintained, and operated by trained and competent persons.

\paragraph{Key frameworks relevant to this chapter include: CDM 2015; MHSWR 1999; HSG47; BS 6031:2009; PUWER 1998; Work in Compressed Air Regulations 1996.}

\begin{itemize}
  \bitem{CDM 2015}{Requires pre-construction information to identify known underground hazards; requires the Construction Phase Plan to address excavation management arrangements; imposes duties on principal contractors and contractors to plan, manage, monitor, and coordinate excavation work safely.}
  \bitem{HSG47 Avoiding Danger from Underground Services}{Establishes the industry-standard approach for identifying and avoiding buried services: obtaining drawings from utility owners; CAT and Genny scanning; trial holes by hand excavation; defined safe dig zones around known services. Compliance with HSG47 is treated by the HSE as the minimum acceptable standard.}
  \bitem{BS 6031:2009 Code of Practice for Earthworks}{Provides technical guidance on soil classification, slope stability, shoring design, and earthwork execution; defines the technical competency standard that groundworks supervisors and designers should demonstrate in relation to excavation stability.}
  \bitem{PUWER 1998}{Requires that all excavation plant --- mechanical excavators, compactors, shoring equipment, pumping equipment --- is suitable for purpose, maintained in an efficient state, inspected at appropriate intervals, and operated only by persons with appropriate training and competency.}
\end{itemize}

%---------------------------- SECTION ----------------------------
\section{Conducting the assessment using the five-step method}

\subsection{Step 1 --- Identify hazards}
\paragraph{Excavation hazards in construction include: collapse of unsupported or inadequately supported trench sides causing burial by asphyxiation or crush injury; contact with live buried services (electricity, gas, water, telecoms) causing electrocution, explosion, flooding, or pressure injury; falls of persons into excavations due to absent or inadequate edge protection; falls of objects (plant, spoil, materials) into excavations where operatives are working; flooding and groundwater ingress causing destabilisation of shoring and drowning risk; oxygen-deficient or toxic atmospheres in deep or flooded excavations with confined-space characteristics; undermining or collapse of adjacent structures, foundations, or services due to excavation proximity; plant instability near excavation edges due to proximity of wheels or tracks to the unsupported face.}

\subsection{Step 2 --- Decide who or what is affected}
\paragraph{Persons at risk from excavation hazards include: groundworkers and operatives working within or adjacent to the excavation; plant operators whose machinery may be at risk of overturning near unsupported edges; other trades and operatives accessing the site who may fall into unprotected excavations; members of the public adjacent to excavation work in public spaces or highways; utility workers attending to services identified or struck during excavation; and occupants of adjacent premises affected by ground movement or service disruption caused by excavation. Special consideration must be given to the vulnerability of any person who may approach an unprotected excavation after working hours, including members of the public and site security personnel.}

\subsection{Step 3 --- Evaluate likelihood and consequence}
\paragraph{The consequence of excavation collapse at depth is almost invariably fatal or life-changing, placing the inherent risk rating at the maximum level (25 on a 5×5 matrix) for unshored excavations deeper than 1.2 metres. The likelihood of collapse increases significantly with depth, soil type (made ground, granular soils, saturated clays), proximity of surcharge loads, and adverse weather. Striking a buried high-voltage electricity cable must also be treated as having the potential for fatal consequence. The effectiveness of controls is high: properly designed and installed shoring or well-established battering schedules reduce the probability of collapse dramatically. The residual risk following proper shoring, verified service avoidance, and competent supervision can be reduced to Low (4--6 on a 5×5 matrix).}

\subsection{Step 4 --- Select and implement controls}
\paragraph{Controls for excavation must be applied in sequence. Before excavation commences: obtain service drawings from all relevant utility undertakers; conduct a CAT and Genny survey of the excavation route; excavate trial holes by hand where services are indicated; consult a structural engineer where excavations are within 3 metres of existing structures or where excavations exceed 4 metres in depth; prepare a shoring or battering design from a competent person. During excavation: maintain shoring or battering in accordance with the design; ensure spoil heaps and plant are kept at minimum 1.2 metre setback from excavation edges (or greater as specified by the engineer); provide safe means of access and egress every 6 metres; install edge protection barriers around the perimeter; check excavation sides before each shift for signs of movement, cracking, or water ingress; use atmosphere monitoring equipment in excavations with confined-space characteristics.}

\subsection{Step 5 --- Record and review}
\paragraph{The excavation risk assessment must be reviewed at the start of each working day by the site supervisor; any change in ground conditions, water levels, or excavation geometry must trigger an immediate reassessment. Excavations must be inspected at the start of each shift, after any event likely to have affected their stability, and after any period of adverse weather, with inspection records retained on site. Any change in proximity to structures, services, or surcharge loads must trigger an engineering review before work continues. All service drawings, CAT survey records, trial hole records, and inspection sheets must be retained as part of the site records.}

\example{Five-Step Application}{A drainage contractor is installing a surface water drainage system in a new housing development. Step 1 identifies: collapse of trench sides at 2.0 metres depth in loose sandy fill; contact with a water main identified on the utility drawings at 1.1 metres depth; falls of operatives into the open trench. Step 2 identifies: two groundworkers in the trench; the mechanical excavator operator; site manager conducting daily inspection. Step 3 rates inherent collapse risk at Very High (25/25) for unsupported 2.0 metre trench. Step 4 implements: Mabey Hire Trench Box installed before any operative enters the excavation; hand excavation within 500\,mm of the identified water main; CAT scan re-run after each 5-metre extension of the trench; edge protection barriers installed along the full length of the open excavation; safe ladder access every 6 metres. Residual risk rated Low (4/25). Step 5 requires daily inspection records before each shift.}

%---------------------------- SECTION ----------------------------
\section{Applying the hierarchy of controls}

\paragraph{The hierarchy of controls for excavation work must begin with the question of whether the excavation can be designed or redesigned to eliminate the deepest or most hazardous sections; for example, directional boring or pipe-bursting techniques can avoid the need for open-cut excavation of deep sections adjacent to structures.}

\begin{itemize}
  \bitem{Elimination}{Avoid open excavation entirely where technically feasible: specify directional boring, horizontal directional drilling, or pipe-bursting for service installation beneath existing structures or roads; use driven or jacked pile systems to avoid deep open excavations for foundations where ground conditions permit.}
  \bitem{Substitution}{Replace deep open-cut excavation with shallower routes where the service alignment allows; replace unsupported battering of trench sides with hydraulic shoring box systems where battering setback requirements would encroach on adjacent features; replace manual excavation near live services with vacuum excavation systems to eliminate strike risk.}
  \bitem{Engineering controls}{Install trench boxes, hydraulic shoring, or ground-anchored sheet piling before any operative enters the excavation; batter trench sides in accordance with soil classification and BS 6031 schedules where space permits; dewater excavations to remove standing water and reduce hydrostatic pressure on trench sides; install edge protection barriers at all exposed excavation perimeters; provide safe and sufficient means of egress every 6 metres.}
  \bitem{Administrative controls}{Obtain and verify service drawings from all utility undertakers before excavation commences; conduct CAT and Genny survey of the entire excavation route; excavate trial holes by hand in all areas of confirmed or suspected service presence; conduct pre-shift inspections and record findings before any operative enters the excavation; restrict access to excavations using barriers and signage; hold permit-to-dig for all excavations in locations with known or suspected buried services.}
  \bitem{PPE}{Hard hat, high-visibility vest, and steel toe-capped boots as minimum personal protective equipment for all groundwork operatives; hearing protection where mechanical excavation equipment is operating at close quarters; atmosphere monitoring equipment worn or carried by any operative entering a deep excavation with confined-space characteristics; appropriate PPE for contact with contaminated ground where the site history indicates contamination risk.}
\end{itemize}

%---------------------------- SECTION ----------------------------
\section{Cadence of assessment and review triggers}

\paragraph{Excavation conditions can change significantly within a single working day as ground conditions respond to weather, dewatering, adjacent working, or changes in the excavation geometry. The assessment review cadence must reflect this dynamic environment.}

\begin{itemize}
  \bitem{Pre-shift daily inspection}{Before any operative enters the excavation at the start of each shift, the site supervisor must inspect the condition of trench sides, shoring, edge protection, access provision, and groundwater levels, and confirm that conditions match those described in the risk assessment. The inspection must be recorded.}
  \bitem{After adverse weather}{Following any significant rainfall, frost, or thaw, the excavation must be re-inspected before work resumes; saturation can cause previously stable cohesive soils to lose shear strength rapidly, and freeze-thaw cycles can destabilise shoring connections and soil structure.}
  \bitem{Change in excavation geometry}{When the excavation is extended, deepened, or widened beyond the scope of the original assessment, the risk assessment must be reviewed and updated before work on the extension commences; changes in depth or proximity to structures may require an engineering design review.}
  \bitem{Discovery of unrecorded services or ground conditions}{Any discovery of buried services not identified in the pre-excavation survey, or any change in ground conditions (e.g.\ unexpected fill, voids, contamination) must trigger an immediate stop, reassessment, and determination of the appropriate response before excavation resumes.}
  \bitem{After any near-miss or incident}{Any near-miss event involving excavation stability, service contact, or falls into the excavation must trigger an immediate review of the risk assessment and must be notified to the principal contractor before work resumes.}
\end{itemize}

%---------------------------- SECTION ----------------------------
\section{Common failure modes}

\begin{itemize}
  \bitem{Unsupported trenches entered without shoring}{The most common and consistently fatal failure in excavation work: operatives enter trenches deeper than 1.2 metres without shoring or battering, relying on apparent soil stability. The risk assessment must make shoring or battering a hard requirement for excavations deeper than 1.2 metres, without exception for ``short duration'' entry.}
  \bitem{Service drawings not obtained or relied upon exclusively}{Failure to obtain updated service drawings from all relevant undertakers, or relying on drawings alone without a CAT and Genny scan, is the primary cause of buried service contact. Service drawings are rarely current; they must be treated as indicative only and verified by scanning and, where services are indicated, trial holes.}
  \bitem{Spoil heap deposited adjacent to the trench edge}{Spoil heaps and plant positioned too close to the trench edge add surcharge loading that can precipitate collapse. The risk assessment must specify a minimum setback distance for spoil and plant based on soil type and shoring design.}
  \bitem{Absence of safe means of egress}{Operatives working in excavations without a ladder or step within reach cannot escape rapidly in the event of ground movement or water ingress. The provision of means of egress every 6 metres is a basic requirement that is routinely neglected on short-duration excavations.}
  \bitem{Atmosphere not monitored in deep or confined excavations}{Excavations more than 1.2 metres deep that are adjacent to sewers, in contaminated ground, or subject to flooding can accumulate oxygen-deficient or toxic atmospheres. Failure to test the atmosphere before entry and to monitor it continuously during work is a common and potentially fatal oversight.}
  \bitem{Pre-shift inspection not conducted or not recorded}{An inspection that is not recorded provides no evidence that it was conducted and creates no audit trail for the condition of the excavation over time. Inspections that are conducted but not recorded are treated as not conducted for enforcement purposes.}
\end{itemize}

\example{Failure Pattern}{A utilities contractor is relaying gas mains in a town centre street. The gas main is specified at 750\,mm depth, but localised undulation of the carriageway has resulted in the main being at 550\,mm in one section. The contractor's risk assessment specifies hand excavation within 500\,mm of any gas main, but the operatives interpret the 750\,mm datum as absolute and switch to mechanical excavation at 250\,mm depth. The mechanical bucket strikes the main at 550\,mm, causing a gas release. The street is evacuated, an emergency gas engineer is called, and the road is closed for six hours. The risk assessment failure is the assumption that the specified depth of a service is invariant and reliable --- an assumption explicitly contradicted by HSG47.}

%---------------------------- CHAPTER SUMMARY ----------------------------
\section{Chapter Summary}

\begin{table}[H]
\centering
\begin{tabularx}{\textwidth}{>{\raggedright\arraybackslash}p{3.2cm}X}
\toprule
\textbf{Assessment Element} & \textbf{Operational Requirement} \\
\midrule
Legal/Professional Basis & CDM 2015; MHSWR 1999; HSG47; BS 6031:2009; PUWER 1998. All excavations must be assessed before commencement; shoring or battering required for depths exceeding 1.2 metres. \\
Method & Apply the full five-step process and document inherent/residual risk. \\
Controls & Shoring or battering to design; service drawings, CAT scan, and trial holes before excavation; safe egress every 6 metres; edge protection barriers; daily pre-shift inspection. \\
Cadence & Pre-shift daily inspection and record; review after adverse weather, change in geometry, discovery of unrecorded services, or any near-miss. \\
Assurance & Use audit, incident learning, and governance reporting to verify effectiveness. \\
\bottomrule
\end{tabularx}
\end{table}

\begin{itemize}
  \bitem{Key takeaway 1}{Unsupported excavations deeper than 1.2 metres must be treated as a near-certain-fatality scenario; shoring or battering is a non-negotiable control, not a cost variable.}
  \bitem{Key takeaway 2}{Service drawings from utility undertakers must be treated as indicative only; CAT and Genny scanning and trial holes by hand excavation are required before any mechanical excavation adjacent to indicated services.}
  \bitem{Key takeaway 3}{Daily pre-shift inspection and recording is a regulatory and practical requirement; ground conditions change overnight, and operatives must not enter excavations until inspection confirms the previous day's conditions are unchanged.}
  \bitem{Key takeaway 4}{Spoil heaps and plant must be kept at minimum 1.2 metres from excavation edges; surcharge loading from adjacent equipment and materials is a significant and preventable cause of trench collapse.}
\end{itemize}

\barquote{In Chapter~\ref{chap:Manual Handling Risk Assessment}, we turn from this assessment to the next domain in the Prudentia framework.}
